%------------------------------------------------------------------------- 
\Section{Language}
\label{section:language}

We analyze concurrent programs written in a stripped version of the C
language and using POSIX threads.
%
% Figure~\ref{fig:SMLsyntax} describes
% this SML language (Simple Multithreaded Language).
We call it here the SML language (Simple Multithreaded Language).
%
We use a relatively small set of operations which follows closely
those of an instruction set architecture (ISA). Intuitively, we filter
out the C language and Pthreads constructs. An ISA includes arithmetic
(e.g.\ \emph{add} and \emph{sub}) and logic instructions (e.g.\
\emph{and}, \emph{or}, \emph{not}), data instructions (e.g.\
\emph{load}, \emph{store}) and control flow instructions (e.g.\
\emph{goto}, \emph{branch}).
%
Arithmetic and logic instructions are simply pure CPU operations. We
are interested in instructions touching the main memory and
instructions controlling the flow. Hence, the language we allow
abstracts away the CPU operations and narrows down the operations to a
set of \emph{movers}, i.e.\ loading from and storing to the main
memory and operations related to controlling the flow, thread
synchronization and thread bookkeeping.

% \begin{figure}
%   \begin{tikzpicture}
%     \node[text width=0.95\linewidth,rounded corners,inner ysep=1pt,inner xsep=1pt,bottom color=blue!20,top color=gray!10!white]
%     {\lstinputlisting{examples/grammar.txt}};
%   \end{tikzpicture}
% \caption{Simple Multithreaded Language}
% \label{fig:SMLsyntax}
% \end{figure}

%% ---------------
We do not define formally the semantics of SML programs. They allow
multiple threads to execute concurrently and manipulate two kinds of
variables: local and global. A global variable is
shared %and can be accessed
by all threads. A local variable is local to a given thread and cannot
be accessed by other threads.
%
%% ---------------
%% Read and write
A shared variable $x$ is a global variable that can be read and
written by any thread. We abstract away the data and extract the mover
instructions as \emph{Read x} or \emph{Write x} accordingly. Local
variables are useful for the control flow only and therefore
abstracted away.
%% ---------------
%% Branching
The \emph{if}, \emph{if-then-else} and \emph{while} statements are
allowed in the form of \emph{branch} and \emph{goto}
combinations. \emph{for}-loops can be unrolled and equivalently
implemented with \emph{while} and \emph{goto}.

%% ---------------
%% Locks
Threads can use locks for synchronization purposes through
\emph{acquire} and \emph{release} primitives. A lock (also called
\emph{mutex}) is a special shared variable that has two values: it is
either \emph{free} or \emph{busy}. A thread trying to acquire a lock
will block if the lock is busy. It will acquire the lock (i.e.\
atomically set it to busy) if the lock was free and continue
execution. Releasing the lock resets it to free.

%% ---------------
%% Condition variable
Threads can use condition variables for synchronization purposes
through \emph{wait} and \emph{signal} primitives.
%
According to the POSIX semantics, wait shall be called while holding a
lock or undefined behavior results.  If a thread waits on a condition
variable, it releases the lock and blocks its execution. A thread will
try to re-acquire the lock (leading to eventual delay) and resume
execution if the condition variable is signaled by another thread.
%
We do not allow at the moment to broadcast a wake up signal (i.e.\
wake up all threads waiting on the condition variable). Signaling a
condition variable, while no thread is waiting on that condition
variable, has no effect.
%
The wait instruction suffers from spurious wakeups as it may return
when no thread specifically signalled that condition variable. We do
not disallow this behaviour and will handle it in our model. It simply
forces the programmer to check a predicate invariant %(in a loop)
upon return and it is not orthogonal to our analysis.

%% ---------------
%% Thread creation
Finally, a thread can create another thread and continue its
execution.
%
% A thread can \emph{join} another thread, in which case it suspends
% execution and resumes it when the other thread terminates (normally)
% or when explicitely calling \emph{exit}. We do not allow at the
% moment to wait for termination of multiple threads.
%
We do not allow at the moment to wait for termination of newly created
threads.

%% ---------------
%% Examples

Figure~\ref{fig:languageExampleCS} and
Figure~\ref{fig:languageExampleProdsCons} show examples written in C
using Pthreads and its equivalent using the SML language. We do not
show the whole code but just the part of interest.
%
Figure~\ref{fig:languageExampleCS} shows a critical section problem
where multiple threads try to update a counter. Only one thread is
allowed to update the counter at a time, to avoid a data race.
%
Figure~\ref{fig:languageExampleProdsCons} is a producers/consumers
example. The producers update a shared buffer, if the buffer is empty
and the consumer reads the buffer if the buffer is full.


\begin{figure}
  \begin{tikzpicture}
    [every node/.append style={rounded corners,inner sep=1pt,bottom color=blue!20,top color=gray!10!white}]
    \node[text width=0.45\linewidth,below left] (cs) {\lstinputlisting{examples/cs-pth.c}};
    \node[text width=0.45\linewidth,below right] (cs2) at(0.5,0){\lstinputlisting{examples/cs.c}};
  \end{tikzpicture}
  \caption{Critical section problem in pthreaded code (left) and its SML counterpart (right).}
  \label{fig:languageExampleCS}
\end{figure}

\begin{figure}
  \begin{tikzpicture}[every node/.append style={inner sep=0pt}]
    \begin{scope}
      \node[text width=0.55\linewidth,above,font=\scriptsize] (conds) {pthread\_cond\_t cvEmpty, cvFull; };
      \node[text width=0.55\linewidth,above,font=\scriptsize] (shared) at (conds.north){int buffer; pthread\_mutex\_t L; };
      \node[text width=0.55\linewidth,below] (prods) {\lstinputlisting{examples/prods-pth.c}};
      \node[text width=0.55\linewidth,below] (cons) at (prods.south){\lstinputlisting{examples/cons-pth.c}};
    \end{scope}
    
    \begin{scope}[xshift=0.52\linewidth]
      \node[text width=0.44\linewidth,above,font=\scriptsize] (shared2) {shared buffer, L, cvEmpty, cvFull; };
      \node[text width=0.44\linewidth,below] (prods2) {\lstinputlisting{examples/prods.c}};
      \node[text width=0.44\linewidth,below] (cons2) at (prods2.south) {\lstinputlisting{examples/cons.c}};
    \end{scope}

    \begin{pgfonlayer}{my background}
      \node[fit=(conds) (shared) (prods) (cons),rounded corners,inner ysep=1pt,inner xsep=1pt,bottom color=blue!20,top color=gray!10!white] (r){};
      \node[fit= (shared2) (prods2) (cons2),rounded corners,inner ysep=1pt,inner xsep=1pt,bottom color=blue!20,top color=gray!10!white] (r2){};
    \end{pgfonlayer}
    
  \end{tikzpicture}
  \caption{Producers/Consumers in pthreaded code (left) and its SML counterpart (right).}
  \label{fig:languageExampleProdsCons}
\end{figure}
